% SPDX-License-Identifier: Apache-2.0
\section{A Map Said Once}
\label{sec:x-regmap}

Every peripheral on AXI-Lite in this document decodes its registers
the same way: a \code{select!} over the address's word bits for the
read, a guard of \code{wsel == k} for every register a write reaches,
and the map stated again in the comment above, in the datasheet, in
the driver and in the device tree. Eleven of them wrote it by hand,
and one of them answered the wrong word for an address it did not
name, which nobody read off the source: the comment said zero and the
hardware said otherwise.

\code{regmap!} is the map said once. A declaration names the map and
its two functions, says how many address bits select a word, and
lists the registers: a word index, a name, an access and a sentence,
and under a register its fields, each a name, a low bit, a width, an
access, a reset value and a sentence.
From that one text the macro writes a module of offsets a program
addresses by, \code{knobs::ctrl} and the rest; the read mux,
\code{knobs\_read}, which takes the registers in declaration order
and answers zero for a word not named; the write enables,
\code{knobs\_we}, a bit a register that \code{with!} guards on; a
constant a field, \code{knobs::ctrl\_step}, with its shift, its width
and its reset, which gets a field out of a word and sets it in one; a
function a field, \code{knobs\_ctrl\_step}, the field out of a written
word, so a unit updates one field under the register's enable and
never reads to modify; a function a register with fields,
\code{knobs\_ctrl\_pack}, the fields packed into the word the read mux
answers; and a table, \code{knobs::MAP}, that the tools read. The lowering reads the
same declaration and inlines the same mux and enables into the
netlist, so the map a program sees and the map the hardware decodes
cannot differ. \code{//tools/regmap} writes the table as a C header
for a driver and as a device tree node at a base, and
\code{//tools/datasheet} writes it as the sheet's register table.

The example is the smallest peripheral that has a map: an identity, a
control word of two fields, a run bit and a step, and a counter that
adds the step while the bit is set. Its \code{run} decodes nothing by
hand. The host client reads the identity, writes the two fields with
the constants' \code{with}, reads the word back and takes them apart
with \code{get}, reads the count twice, stops the counter, and reads
a word the map does not name; the run prints what it saw, then the
header, a define a register and four a field, and the table rows the
declaration stands for.

\begin{figure}[htbp]
\centering
\input{regmap_timing.tex}
\caption{The control word written on \code{aw}, its two fields
landing in \code{run} and \code{step}, the counter running by the
step, and the reads on \code{ar} answered on \code{r}.}
\label{fig:x-regmap}
\end{figure}

\lstinputlisting[caption={\code{ex\_regmap.rs}. Run with \code{bazel run //lib/examples:ex\_regmap}.},label={lst:x-regmap}]{lib/examples/ex_regmap.rs}

\lstinputlisting[style=ascii,caption={What \code{ex\_regmap} printed when the build ran it: the run, then the C header with its fields and the table rows written from the declaration.},label={lst:x-regmap-out}]{out_regmap.txt}
